% Preserved from main.tex at 6a9df36; excluded from the technical-review draft.
\section{A remark on designed topologies}
\label{app:topology}

Nothing in this appendix is about the Lightning Network as it is. The public channel graph is not cubic, not bipartite, and not planar, and the theorem recalled below says nothing about it. The remark concerns an overlay that someone designs, such as the channel set an Ark server places (Section~\ref{sec:ark}) or a factory whose members a warden admits (Section~\ref{sec:bip448}), where the shape of the graph is a choice. It records one property such a shape can be given, what the property amounts to in the terms of this paper, and what it costs. No proposition in the main text depends on it.

\paragraph{The result and its standing.} A preprint by OpenAI \citep{openai2026barnette} claims Barnette's conjecture: every finite simple cubic bipartite planar 3-vertex-connected graph has a Hamiltonian cycle (its Theorem~1.1). Its Corollary~7.1 adds that for every edge there is a Hamiltonian cycle avoiding that edge. The manuscript is an unrefereed, machine-generated preprint. We checked the statements of Theorem~1.1 and Corollary~7.1 against the manuscript and read its outline; we have not refereed the proof. A Lean formalization is published beside it: thirty-one modules in which a text scan finds no \texttt{sorry}, \texttt{admit}, or added axiom, reached from a challenge file that states the theorem, leaves its own proof as \texttt{sorry}, and names those modules as the solution. We compiled neither. No outside review is known to us. The paper says of its argument that it ``does not assert a polynomial-time construction of the cycle.'' Everything below is therefore conditional on a claim whose checking is incomplete, and is written so that nothing else in the paper leans on it.

\paragraph{What it would give a designed overlay.} Let the overlay be a channel graph in which every participant holds exactly three channels, the participants fall into two classes with every channel between the classes, the graph is planar, and no two participants' removal disconnects it. If Theorem~1.1 holds, there is a closed route that visits every participant once. In the notation of Appendix~\ref{app:notation}, a payment of $x$ from each participant to its successor along the route is a settled event on each channel of the route and on no other. Every participant's net position is unchanged, since each pays $x$ and receives $x$; on each route channel $x$ moves from the upstream to the downstream side; and each participant's spendable balance is thereby rotated from its successor channel to its predecessor channel. If Corollary~7.1 holds, such a route exists that avoids any one channel named in advance, so after any single channel closes or is left out of the design some route remains, though not the same one. That is the whole of what the theorem supplies: the existence of one spanning route for a uniform circulation. It does not supply a route through a channel the designer names, and it does not supply the several routes that an arbitrary pattern of imbalances would need, since one circulation shifts every route channel by the same amount in the same rotational sense. The paper's Corollary~7.2, which places every three-edge path on some Hamiltonian cycle for a class that includes the planar case, would give routes through a named channel; it rests on an equivalence the paper takes from elsewhere, and we do not use it.

\paragraph{Limits.} Three. Existence is not construction: the argument asserts no polynomial-time construction, so a designer who wants the route has to search for it, and on an overlay of any size the theorem promises only that the search will end. The shape is expensive in the terms of Section~\ref{sec:lightning}: every participant holds three two-party channels, each direction-partitioned, so the forecast of Section~\ref{sec:lightning} is made three times per participant where the star of Section~\ref{sec:ark} makes it once and the hyperedge of Section~\ref{sec:bip448} once per group; a cubic bipartite graph has classes of equal size, so if one class is agents and the other supplies their liquidity, the design needs as many suppliers as agents; and nothing in it pools, so Remark~\ref{rem:sqrtN} offers it nothing. The route is a property the public graph cannot be given, bought at roughly three times the per-agent channel's capital per participant. And a route that exists as a cycle of the graph need not be feasible at a given amount. Each channel on it has to hold at least $x$ on its upstream side when the payment is made, which is the condition of Lemma~\ref{lem:cut} on each edge of the ring, and a circulation feasible in one rotational sense may not be in the other. \citet{dugan2026trial} makes the general point against the earlier manuscript's cut reasoning: feasibility depends on every prefix of the arrival sequence and on which edge the liquidity sits, so that two histories with the same totals and the same final imbalance can differ in what they serve, and a bound on a cut is not a certificate of delivery. The theorem speaks to the graph and not to the balances on it, and the balances are what Sections~\ref{sec:lightning} to~\ref{sec:bip448} are about.

\paragraph{Witness.} The reference model checks the statement on three small graphs of the class, the cube, the hexagonal prism, and the truncated octahedron: for each edge in turn it finds a Hamiltonian cycle avoiding that edge by exhaustive search, and checks that a unit circulation round the cycle nets to zero at every participant, moves one unit across each cycle channel, and moves nothing elsewhere (Appendix~\ref{app:model}). It witnesses the statement on finite cases and proves nothing further.

